\specialchapter{记号索引}

\(\mathcal{L}^{\mathrm{c}}(\mathrm{E};\mathrm{F})\) (紧线性映射) 2 \(\mathrm{E}_{(\mathbf{C})}\) (复化向量空间) 2 \(u_{(\mathbf{C})}\) (复化线性映射) 3 \(\mathcal{L}^{\mathrm{f}}(\mathrm{E};\mathrm{F})\) (有限秩线性映射) 4 \(\mathcal{C}_0(\mathrm{X};\mathrm{K}), \mathcal{C}_0(\mathrm{X})\) (在无穷远处趋于0的连续函数) 20 \(\mathcal{N}^{p,q}(\mathrm{X} \times \mathrm{Y})\) (核空间) 27 \(u_k\) (带核的积分算子) 29 \(k_\circ\)（映射 \(y \mapsto k_y\) ，其中 \(k_y(x) = k(x,y)\) ）30 \(\mathcal{L}^{p',q'}(\mathrm{X},\mathrm{Y},\mu,\nu)\) (核空间) 31 \(u \equiv v\) (模 \(\mathcal{L}^{\mathrm{f}}(\mathrm{E};\mathrm{F})\) 的同余) 40 \(\mathcal{F}(\mathrm{E};\mathrm{F})\) (弗雷德霍姆映射) 41\\
ind(u) (弗雷德霍姆映射的指标) 43 \(\mathcal{I}(\mathrm{E};\mathrm{F})\) (E 到 F 的同构) 57 \(\mathcal{M}\mathcal{D}(\mathrm{E};\mathrm{F})\) (单射直态射) 57 \(\mathcal{E}\mathcal{D}(\mathrm{E};\mathrm{F})\) (满射直态射) 57 \(\mathcal{R}(\mathrm{E})\) (里斯自同态) 60\\
((u)) (u 的余范数) 61 \(\mathcal{M}(\mathrm{E};\mathrm{F})\) (严格单射态射) 67 \(\mathcal{Q}\mathcal{M}(\mathrm{E};\mathrm{F})\) (核有限维的严格态射) 67 \(\mathcal{E}(\mathrm{E};\mathrm{F})\) (连续满射态射) 70 \(\mathcal{Q}\mathcal{E}(\mathrm{E};\mathrm{F})\) (像具有有限余维的态射) 70 \(\mathcal{C}alk(\mathrm{E})\) (巴拿赫空间 E 的卡尔金代数) 75 \(\mathrm{N}_{\lambda}\) (幂零空间) 78 \(\mathrm{I}_{\lambda}\) (余幂零空间) 78 \(e_{\lambda}(u)\) (谱投影) 82 \(m_{\lambda}(u)\) (谱重数) 82 \(\mathrm{Sp}_{\mathrm{s}}(u)\) (谱的敏感点) 83 \(\mathrm{Sp}_{\mathrm{e}}(u)\) (本质谱) 83 \(\mathrm{Sp}_{n}(u)\) (谱中使得 \(u - \lambda 1_{\mathrm{E}}\) 具有指标 n 的元素) 86 \(\mathrm{Sp}_{\omega}(u)\) (\(\mathrm{Sp}_{n}(u)\) 的并集的补集) 86 \(\mathcal{D}_{\mathrm{B}}(\mathrm{E})\) (基 B 中对角的自同态) 146 \(\mathrm{I}_{\mathrm{E}}\) (F ≠ E 的有限维子空间的维数) 151 \(\lambda_n(u)\) (u 的特征值递减序列) 152

\(r_{\mathrm{F}}(u), \mathrm{R}_{\mathrm{F}}(u)\) （瑞利商）\ldots.. 153 \(\mathcal{F}_n\) （维数 \(n\) 的子空间）\ldots.. 153 \((\alpha_n(u))_{n \in \mathrm{J}}\) （奇异值序列）\ldots.. 156 \((\alpha_n(u))_{n \in \mathrm{I_E}}\) （奇异值的扩展序列）\ldots.. 156 \(\mathcal{L}_1(\mathrm{E})\) （有限迹自同态）\ldots.. 164 \(\|u\|_1\) （\(\mathcal{L}_1(\mathrm{E};\mathrm{F})\) 上的范数）\ldots.. 167 \(\mathcal{L}_1(\mathrm{E};\mathrm{F})\) （核型线性映射）\ldots.. 169 \(\mathcal{L}_{\mathrm{u}}(\mathrm{X})\) （普遍可测函数）\ldots.. 183 \(\mathcal{L}_{\mathrm{u}}^{\infty}(\mathrm{X})\) （有界普遍可测函数）\ldots.. 183 \(m_g, \widetilde{m}_g\) （乘法自同态）\ldots.. 186 \(\mathcal{G}_{\mathrm{A}}\) （盖尔范德变换）\ldots.. 190 \(\mathcal{H}_{\mathrm{A}}\) （盖尔范德变换的逆）\ldots.. 190 \(\mathcal{F}\) （傅里叶变换）\ldots.. 196 \(\overline{\mathcal{F}}\) （余傅里叶变换）\ldots.. 196 \(\mathcal{D}(\mathrm{U})\) （U 上的测试函数）\ldots.. 200 \(\mathcal{D}_{\mathbf{R}}(\mathrm{U})\) （实值测试函数）\ldots.. 201 \(\mathcal{D}'(\mathrm{U})\) （U 上的广义函数）\ldots.. 203 \(\partial^\alpha f\) （广义函数的偏导数）\ldots.. 208 \(\mathcal{S}(\mathbf{R}^n)\) （\(\mathbf{R}^n\) 上的施瓦茨函数）\ldots.. 209 \(\mathcal{S}'(\mathbf{R}^n)\) （\(\mathbf{R}^n\) 上的缓增广义函数）\ldots.. 215 \(\mathcal{M}^t(\mathbf{R}^n)\) （缓增测度空间）\ldots.. 216 \(W^{k,p}(U)\) （指标 \(k\) 、指数 \(p\) 的索伯列夫空间）\ldots.. 222 \(W_0^{k,p}(U)\) （\(\mathcal{D}(U)\) 在 \(W^{k,p}(U)\) 中的闭包）\ldots.. 222 \(H^k(U)\) （指标 \(k\) 、指数2的索伯列夫空间）\ldots.. 222 \(H_0^k(U)\) （\(\mathcal{D}(U)\) 在 \(H^k(U)\) 中的闭包）\ldots.. 222 \(\mathcal{P}(\mathrm{E};\mathrm{F})\) （从 E 到 F 的部分算子）\ldots.. 224 \(\mathcal{P}(\mathrm{E})\) （E 上的部分算子）\ldots.. 224 dom( \(u\) ) （\(u\) 的定义域）\ldots.. 224 \(1_{\mathrm{D}}\) （定义域为 D 的恒等部分算子）\ldots.. 225 \(u \circ v\) （部分算子 \(u\) 与 \(v\) 的复合）\ldots.. 225 \(u \subset v\) （\(v\) 是 \(u\) 的扩张）\ldots.. 226 \((u_i)_{i \in I}\) （乘积部分算子）\ldots.. 226 \(u + v\) （部分算子的和）\(u\)、\(v\)\ldots.. 226 \(\overline{u}\) （\(u\) 的闭包）\ldots.. 228 \(E_u\) （由 \(u\) 的定义域定义的赋范空间）\ldots.. 230 \(\|x\|_u\) （空间 \(E_u\) 上的范数）\ldots.. 230 \((x | y)_u\) （定义希尔伯特空间 \(E_u\) 的内积）\ldots.. 230 \(m_g, \widetilde{m}_g\) （乘法部分算子）\ldots.. 231 \(u^*\) （\(u\) 的伴随）\ldots.. 235 \(\mathcal{A}(\mathrm{E})\) （自伴部分算子）\ldots.. 238 \(^t\mathrm{D}\) （转置微分算子）\ldots.. 242 \(H_D\) （微分算子 \(D_+\) 的定义域）\ldots.. 242 \(D_-\) （定义域为 \(\mathcal{D}(U)\) 的微分算子）\ldots.. 242 \(D_+\) （定义域为 \(H_D\) 的微分算子）\ldots.. 242 \(R(u,\lambda)\) （\(u\) 的预解式）\ldots.. 244 \(\mathrm{PSp}_{\varepsilon}(u)\) （\(\varepsilon\)-伪谱）\ldots.. 250 \(b(u)\) （\(u\) 的有界化）\ldots.. 262 \(\Omega(\mathrm{E})\) （使 \(v\) 正且单射的 \(1_{\mathrm{E}} - v^*v\) 的集合）\ldots.. 264 \(B(v)\) （部分算子 \(v\circ((1_{\mathrm{E}} - v^*v)^{1/2})^{-1}\) ）\ldots.. 264

\(\beta (z) = z / \sqrt{1 + |z|^2}\) 266 \(D_f\) （\(f(u)\) 的定义域）271 \(f(u)\) （普遍可测函数演算）272 \(p_A\) （由 A 定义的谱投影）277 \(u^\alpha\) （函数演算）289 \(|u|\) （\(u\) 的绝对值）290 \(q_u\) （与 \(u\) 关联的正部分形式）291 \(E_q\) （预希尔伯特空间）292 \((x|y)_q\) （\(E_q\) 上的内积）292 \(\| x \| _q\) （\(E_q\) 上的范数）292 \(u_q\) （表示 \(q\) 的部分算子）292 \(\mathrm{Sp}_{\mathrm{s}}(u)\) （\(u\) 的敏感谱）297 \(\mathrm{Sp}_{\mathrm{e}}(u)\) （\(u\) 的本质谱）297 \(\mathrm{Sp_h}(u)\) （谱的上部）298 \(\varrho_{e}\) （本质谱的下界）298 \(\mathrm{Sp_b}(u)\) （谱的下部）298 \(\mathcal{F}_n\) （维数 \(n\) 的子空间）299 \(\mathcal{F}_n^u\) （\(\mathrm{dom}(u)\) 中维数 \(n\) 的子空间）299 \(\widetilde{r}_{\mathrm{F}}(u), \widetilde{\mathrm{R}}_{\mathrm{F}}(u)\) （瑞利商）300 \(\|v\|_u\) （关于 \(v\) 的 \(u\) 的相对范数）306\\
vp （柯西主值）337 \(\chi_\varrho\) （表示的特征标）374 \(\mathrm{Res}_{H}^{G}(\varrho)\) （\(\varrho\) 到 H 的限制）375 \(\mathrm{Hom}_{G}(\varrho_1,\varrho_2)\) （G-态射空间）375 \(\varrho^n = n \varrho\) （表示 \(\varrho\) 的副本之和）376 \(\check{\varrho}\) （逆步表示）377 \(\mathrm{Sesq_G}(\varrho, \pi), \mathrm{Sesq_G(E,F)}\) （G-不变半双线性形式）380 \(\overline{\varrho}\) （共轭表示）382 \(\varrho_1 \boxtimes \varrho_2\) （外张量积）385 \(\varrho_1 \otimes \varrho_2\) （张量积）385 \(\varrho^{\otimes n}\) （\(\varrho\) 的张量幂）385 \(\Upsilon(G)\) （矩阵系数）386 \(\Theta(G)\) （有限维矩阵系数）386 \(\operatorname{cl}(\pi)\) （酉表示的类）392 \(\widehat{G}\) （酉对偶）393 \(M_\pi (\varrho)\) （同型分量）394 \(\mathcal{M}_c(G)\) （紧支撑测度）400 \(\mathcal{M}^1(G)\) （有界测度）401 \(\gamma_{G}^{(p)}\) （左正则表示）405 \(\delta_{G}^{(p)}\) （右正则表示）405 \(\varrho_{G}^{(p)}\) （双正则表示）406 \(\gamma_{G}\) （\(L^2\) 中的左正则表示）407 \(\delta_{G}\) （\(L^2\) 中的右正则表示）407 \(\varrho_{G}\) （\(L^2\) 中的双正则表示）407 \(\mathcal{F}_{\pi}(G)\) （π-等变函数）408 \(\mathcal{K}_{\pi}(G)\) （在）\(\mathcal{F}_{\pi}(G)\) 中模 H 具有紧支集的连续函数）\ldots{} 408 \(N_p(f)\) （在）\(\mathcal{F}_π(G)\) 上的半范数）\ldots{} 408 \(\mathcal{L}_{\pi}^{p}(G)\) （模 H 对幂 \(p^e\) 可积的 π-等变函数）\ldots{} 409

\(L_{\pi}^{p}(G)\) （与 \(\mathcal{L}_{\pi}^{p}(G)\) 关联的巴拿赫空间）\ldots.. 410 \(\mathcal{L}_{\pi}^{2}(G)\) （模 H 平方可积的 \(\pi\)-等变函数）\ldots.. 415 \(L_{\pi}^{2}(G)\) （与 \(\mathcal{L}_{\pi}^{2}(G)\) 关联的希尔伯特空间）\ldots.. 415 \(\operatorname{Ind}_{H}^{G}(\pi)\) （由 H 的表示 π 诱导的 G 表示）\ldots.. 416 \(\gamma_{G,\chi}\) （在 \(\mathcal{F}_{\chi}(G)\) 或 \(L_{\chi}^{2}(G)\) 上的左正则表示）\ldots.. 417 \(\delta_{G,\chi}\) （在 \(\mathcal{F}_{\chi}(G)\) 或 \(L_{\chi}^{2}(G)\) 上的右正则表示）\ldots.. 417 \(\gamma_{G,\chi}^{(p)}\) （\(L_{\chi}^{p}(G)\) 中的左正则表示）\ldots.. 417 \(\delta_{G,\chi}^{(p)}\) （\(L_{\chi}^{p}(G)\) 中的右正则表示）\ldots.. 417 \(\gamma_{G,\chi}\) （\(L_{\chi}^{2}(G)\) 中的酉左正则表示）\ldots.. 417 \(\delta_{G,\chi}\) （\(L_{\chi}^{2}(G)\) 中的酉右正则表示）\ldots.. 417 \(\varrho_{G,\chi}\) （\(L_{\chi}^{2}(G)\) 中的酉双正则表示）\ldots.. 417 \(c_{Z}(\pi)\) （形式次数）\ldots.. 423 \(\mathrm{Noy}_{+}(\mathrm{X})\) （普遍正核）\ldots.. 434 \(A_{+}^{\prime}\) （连续正线性泛函）\ldots.. 436 \(\mathrm{Pos}(G)\) （正定函数）\ldots.. 443 \(\mathrm{Pos}_{1}(G)\) （满足 \(\varphi(e)=1\) 的正定函数）\ldots.. 443 \(\varrho_{G}(\pi)\) （\((\overline{\pi}\boxtimes\pi)\)-同型分量的 \(\varrho_{G}\)）\ldots.. 462 \(\Theta(G^{\sharp})\)，\(Z\Theta(G)\) （中心矩阵系数）\ldots.. 467 \(R(G)\) （格罗滕迪克环）\ldots.. 469 \(\mathrm{F}(\widehat{G})\) （End(Eπ) 的乘积代数）\ldots.. 470 \(\mathrm{F}_{b}(\widehat{G})\) （有界族的子代数）\ldots.. 470 \(\|u\|_{\mathrm{F}_{b}}\) （\(\mathrm{F}_{b}(\widehat{G})\) 中的范数）\ldots.. 470 \(\mathrm{F}_{0}(\widehat{G})\) （趋于0的族的子代数）\ldots.. 470 \(\overline{\mathcal{F}}_{G}(\nu)\) （ν 的傅里叶余变换）\ldots.. 471 \(\overline{\mathcal{F}}_{G}(f)\) （f 的傅里叶余变换）\ldots.. 471 \(\|u\|_{2}\) （End(Eπ) 中的范数）\ldots.. 471 \(\mathrm{F}^{2}(\widehat{G})\) （End(Eπ) 的希尔伯特直和）\ldots.. 472 \(\mathcal{F}_{G}(f)\) （f 的傅里叶变换）\ldots.. 473 \(\mathcal{C}(\widehat{G})\) （\(\widehat{G}\) 上的函数）\ldots.. 475 \(M^{1}(G^{\sharp})\) （中心有界测度）\ldots.. 475 \(\mathrm{FS}(\pi)\) （弗罗贝尼乌斯--舒尔指标）\ldots.. 476 \(M_{2k}(\varrho)\) （阶为 2k 的绝对矩）\ldots.. 476 \(S^{2}(E)\) （A-模的对称平方）\ldots.. 478 \(\wedge^{2}E\) （A-模的外平方）\ldots.. 478

